\chapter*{Introduction Générale}
\addcontentsline{toc}{chapter}{Introduction Générale}

La transformation numérique constitue aujourd'hui l'un des leviers majeurs de la compétitivité des organisations. Elle s'accompagne d'une volumétrie croissante de données RH et projet, produites par les outils de gestion, les feuilles de suivi et les échanges collaboratifs. Sans dispositif de consolidation, cette richesse informationnelle demeure difficilement exploitable au moment où une décision doit être prise. Les plateformes d'aide à la décision répondent précisément à cet écart, en transformant des données hétérogènes en indicateurs lisibles, en tendances observables et en recommandations susceptibles d'orienter l'action. Elle redéfinit les modes de production, de collaboration et de pilotage, en imposant une disponibilité permanente de l'information et une capacité d'adaptation accrue face à l'incertitude. Dans ce contexte, les outils d'aide à la décision occupent une place centrale : ils ne se limitent plus à restituer des données historiques, mais visent à contextualiser les signaux faibles, à mettre en évidence les écarts par rapport aux objectifs et à accompagner les décideurs dans le choix des actions les plus pertinentes. Pour les directions opérationnelles et les fonctions RH, cette évolution n'est pas un simple confort technologique ; elle répond à un besoin structurel de réactivité, de transparence et de coordination à l'échelle de l'entreprise.

La gestion des talents, des compétences, des charges de travail et du portefeuille de projets illustre particulièrement cette exigence. Les organisations doivent concilier la disponibilité des ressources, le développement des savoir-faire, la tenue des engagements contractuels et la réussite des livrables projet, dans un environnement où les priorités évoluent rapidement. Un retard sur un projet, une montée en charge imprévue ou un déficit de compétences sur une mission critique peuvent avoir des effets en cascade sur la performance collective. Par conséquent, disposer d'une vision consolidée, fiable et actualisée devient un prérequis pour anticiper les tensions, répartir équitablement la charge et aligner les décisions stratégiques avec la réalité du terrain.

Pourtant, les managers et les responsables RH rencontrent encore de nombreuses difficultés dans l'exercice quotidien de ce pilotage. Le manque de visibilité globale demeure une préoccupation récurrente : les informations sont souvent dispersées entre plusieurs outils, feuilles de calcul ou processus informels, ce qui complique la construction d'une lecture commune de la situation. La surcharge de travail constitue un second enjeu, lorsque certaines ressources atteignent des seuils critiques sans alerte précoce ni mécanisme de rééquilibrage. Les fins de contrat, lorsqu'elles ne sont pas anticipées, fragilisent la continuité des projets et accroissent le risque de rupture de compétences. Les écarts entre compétences requises et compétences disponibles freinent le démarrage ou la poursuite des activités, tandis que la difficulté à prioriser les décisions --- arbitrer entre plusieurs demandes RH, choisir un réaménagement d'équipe ou valider un plan de formation --- retarde parfois des actions pourtant nécessaires. Ces contraintes, cumulées, illustrent la nécessité d'une approche intégrée plutôt que d'interventions ponctuelles et déconnectées.

L'intelligence artificielle et l'automatisation ouvrent des perspectives concrètes pour répondre à ces limites. L'analyse assistée des données permet de détecter plus tôt les situations à risque, de formuler des recommandations argumentées et de réduire le temps consacré aux tâches de synthèse répétitives. L'automatisation des workflows, quant à elle, structure les demandes, garantit une meilleure traçabilité des décisions et facilite le passage de l'information entre les acteurs concernés. Il ne s'agit pas de substituer le jugement humain, mais de le soutenir par des indicateurs explicites, des alertes contextualisées et des scénarios d'aide à l'arbitrage, afin que les décisions stratégiques reposent sur des éléments partagés et vérifiables.

C'est dans cette perspective qu'a été conçue la solution développée dans le cadre du présent projet de fin d'études : Strategic Copilot, plateforme SaaS intelligente d'aide à la décision pour la gestion stratégique des talents et des projets. Il s'agit d'un environnement web unifié destiné aux managers, aux équipes RH et aux talents, articulé autour d'espaces métier distincts mais cohérents. Pour les managers, la plateforme propose un pilotage du portefeuille projet, du suivi d'équipe, de la gestion des risques et des demandes adressées aux fonctions RH ou aux talents. Pour les responsables RH, elle couvre le référentiel des talents, le catalogue de compétences, l'identification des écarts critiques, les plans de formation, le traitement des demandes managers et les alertes organisationnelles. Pour les talents, elle offre une visibilité sur les missions, les tâches, la charge de travail, les compétences et les parcours de formation. Au cœur de l'expérience utilisateur, le Strategic Copilot intégré fournit des synthèses, des analyses de viabilité, des recommandations et un assistant conversationnel adapté au contexte métier consulté. Les tableaux de bord, les alertes, le simulateur de scénarios et le journal des décisions complètent ce dispositif en offrant aux utilisateurs une lecture immédiate de la situation, tout en conservant la trace des arbitrages effectués. L'ambition est ainsi de rapprocher, au sein d'un même environnement, la donnée brute, l'interprétation assistée et l'action concrète.

Les objectifs principaux du projet visent à traduire ces orientations en une solution opérationnelle et exploitable. Il s'agit tout d'abord de permettre la visualisation claire des indicateurs RH et projet, afin de disposer d'un tableau de bord partagé par rôle et par périmètre de responsabilité. Il s'agit ensuite de mettre en œuvre une détection proactive des risques, notamment la surcharge, les échéances contractuelles et les manques de compétences, avec des mécanismes d'alerte intelligents. Le projet poursuit également l'analyse de la viabilité des projets, la génération de recommandations automatisées fondées sur l'intelligence artificielle, ainsi que la gestion structurée des demandes RH au sein de workflows décisionnels. Enfin, un objectif transversal consiste à offrir une expérience fluide, sécurisée et évolutive, capable d'accompagner la montée en charge fonctionnelle de la plateforme.

Sur le plan technique, la réalisation s'appuie sur une stack moderne adaptée aux exigences d'une application SaaS. L'interface utilisateur est développée en React.js \cite{react2024}, avec Vite \cite{vite2025} comme outil de build et Tailwind CSS \cite{tailwind2025} pour la conception visuelle et la cohérence des composants. La couche données et services s'appuie sur Supabase \cite{supabase2025}, qui fournit l'authentification, les API et l'accès à une base PostgreSQL \cite{postgresql2025} pour la persistance des informations métier. Les processus d'orchestration et d'intégration sont confiés à n8n \cite{n8n2025}, permettant d'automatiser les flux entre la plateforme, les services externes et les composants d'intelligence artificielle \cite{openai2025,russell2021}. L'enrichissement décisionnel repose sur des modèles de langage invoqués côté orchestrateur, mobilisés pour les recommandations, les synthèses et l'assistant conversationnel du copilot. Cette combinaison technologique a guidé les choix d'architecture retenus tout au long du projet.

La méthodologie générale adoptée suit une démarche progressive, structurée en phases complémentaires. Une première phase d'analyse a permis d'identifier les besoins, les acteurs, les contraintes et les cas d'usage prioritaires associés aux espaces manager, RH et talent. Une phase de conception a ensuite défini l'architecture applicative, les principaux modules fonctionnels, les flux de données et les parcours utilisateurs. La phase de réalisation a été menée de manière itérative, avec des livraisons successives des écrans, des intégrations API et des mécanismes d'alerte et de recommandation. Enfin, des campagnes de tests et de validation ont été conduites sur les scénarios critiques --- pilotage projet, traitement des demandes RH, suivi talent --- afin de vérifier la cohérence fonctionnelle, la stabilité technique et l'apport réel de l'aide à la décision. Cette organisation a favorisé un alignement continu entre les objectifs métier et les choix d'implémentation. Chaque itération a été l'occasion de confronter les prototypes aux attentes des utilisateurs cibles, d'ajuster les écrans de pilotage et de renforcer la fiabilité des échanges entre le front-end, les services Supabase et les flux automatisés orchestrés par n8n. Cette rigueur méthodologique a permis de maintenir une cohérence d'ensemble, malgré la diversité des fonctionnalités à couvrir.

Le présent rapport retrace l'ensemble de cette démarche et en présente les apports de manière structurée. Le premier chapitre est consacré au contexte général du projet et à l'analyse des besoins. Le deuxième chapitre présente l'étude bibliographique : fondements RH, gestion des talents, pilotage de portefeuille, aide à la décision, automatisation et technologies web mobilisées par Strategic Copilot. Le troisième chapitre formalise le cahier des charges, les spécifications et la planification. Le quatrième chapitre présente la conception fonctionnelle et technique complète du système. Le cinquième chapitre détaille la réalisation et l'implémentation : environnement de développement, front-end, workflows n8n, base de données et interfaces livrées. Le sixième chapitre présente les tests réalisés, les résultats obtenus et l'évaluation du système. La conclusion générale revient sur les enseignements du projet, ses limites et les perspectives d'évolution envisageables pour consolider la plateforme comme outil durable de pilotage stratégique des talents et des projets.
